% ---------------------------------------------------------------------------
%  中文学位论文
%  封面、摘要、目录、分章、参考文献（GB/T 7714—2015，biblatex + biber）、致谢。
%  章节在 chapters/，前置与后置部分在 front/，文献条目在 refs.bib。
%  中文字体用系统自带的，西文用内置的 TeX Gyre Termes，不必另装。
%  各校格式要求不同：页边距、字号、页眉在下面的导言区改；封面信息在「论文信息」里改。
%  文中的研究与数据均为示例。
% ---------------------------------------------------------------------------
\documentclass[UTF8,zihao=-4,openany,oneside]{ctexbook}

\usepackage[a4paper,top=3cm,bottom=2.6cm,left=3cm,right=2.6cm,headheight=15pt,headsep=.8cm]{geometry}
\setmainfont{texgyretermes}[
  Extension=.otf, UprightFont=*-regular, ItalicFont=*-italic,
  BoldFont=*-bold, BoldItalicFont=*-bolditalic]
\usepackage{amsmath,amsthm}
\usepackage{unicode-math}
\setmathfont{texgyretermes-math.otf}
\usepackage{xcolor}
\usepackage{graphicx}
\usepackage{booktabs}
\usepackage{caption}
\usepackage{enumitem}
\usepackage{fancyhdr}
\usepackage{pgfplots}
\pgfplotsset{compat=1.18}
\usepackage[backend=biber,style=gb7714-2015]{biblatex}
\addbibresource{refs.bib}
\usepackage[hidelinks]{hyperref}

\graphicspath{{figures/}}
\linespread{1.5}
\setlist{itemsep=0pt, topsep=4pt, leftmargin=2.5em}
\captionsetup{font=small, labelsep=quad}

% 墨色做结构，绿色只用在图表里
\definecolor{ink}{HTML}{1F3A33}
\definecolor{accent}{HTML}{2F7D63}
\definecolor{mist}{HTML}{75837E}

% 章：居中，章号与章名分两行；节、小节左对齐
\ctexset{
  chapter = {
    format     = \centering\color{ink},
    nameformat = \zihao{-4}\color{mist},
    titleformat= \heiti\zihao{3},
    aftername  = \par\vspace{6pt},
    beforeskip = 12pt, afterskip = 30pt,
  },
  section    = {format=\heiti\zihao{4}\color{ink}, beforeskip=18pt, afterskip=10pt},
  subsection = {format=\heiti\zihao{-4}, beforeskip=12pt, afterskip=6pt},
}

% 页眉：左边论文题目，右边当前章；页码在页脚居中
\pagestyle{fancy}
\fancyhf{}
\fancyhead[L]{\small\color{mist}\thesistitle}
\fancyhead[R]{\small\color{mist}\leftmark}
\fancyfoot[C]{\small\thepage}
\renewcommand{\headrule}{\color{mist!50}\hrule height .4pt}
\fancypagestyle{plain}{\fancyhf{}\fancyfoot[C]{\small\thepage}\renewcommand{\headrule}{}}

\theoremstyle{definition}
\newtheorem{definition}{定义}[chapter]
\theoremstyle{plain}
\newtheorem{proposition}{命题}[chapter]

% ---------------------------------------------------------------- 论文信息
\newcommand{\thesistitle}{中文长文本可读性的计算评估方法研究}
\newcommand{\thesistitleen}{Computational Assessment of Readability for Long-form Chinese Text}
\newcommand{\thesisauthor}{苏\,望}
\newcommand{\thesisid}{2023\,0000\,0000}
\newcommand{\thesismajor}{计算机科学与技术}
\newcommand{\thesisadvisor}{何\,明\,远\quad 教授}
\newcommand{\thesisschool}{计算机科学与技术学院}
\newcommand{\thesisuniv}{某\,某\,大\,学}
\newcommand{\thesisdate}{二〇二六年五月}

\begin{document}

\frontmatter
% 封面：内容在 main.tex 的「论文信息」里改，这里只管版式
\begin{titlepage}
  \centering
  \vspace*{1.2cm}
  {\songti\zihao{-2}\color{ink}\thesisuniv\par}
  \vspace{6pt}
  {\zihao{-4}\color{mist}博\,士\,学\,位\,论\,文\par}
  \vspace{2.6cm}

  {\color{ink}\rule{\linewidth}{1.2pt}\par}
  \vspace{18pt}
  {\heiti\zihao{2}\color{ink}\thesistitle\par}
  \vspace{12pt}
  {\zihao{4}\itshape\color{mist}\thesistitleen\par}
  \vspace{18pt}
  {\color{ink}\rule{\linewidth}{.4pt}\par}

  \vfill
  {\zihao{4}
    \renewcommand{\arraystretch}{1.9}
    \begin{tabular}{@{}r@{\hspace{1.5em}}l@{}}
      {\color{mist}作者姓名} & \thesisauthor  \\
      {\color{mist}学\qquad 号} & \thesisid   \\
      {\color{mist}专业名称} & \thesismajor   \\
      {\color{mist}指导教师} & \thesisadvisor \\
      {\color{mist}所在学院} & \thesisschool  \\
    \end{tabular}\par}
  \vfill
  {\zihao{4}\color{ink}\thesisdate\par}
  \vspace*{.6cm}
\end{titlepage}

\chapter{摘\quad 要}

一篇文章读起来累不累，取决于两件事：写了什么，以及排成了什么样子。
已有的可读性公式大多只看前者——句子多长、词有多难——而把版面当作不变的背景。
对于动辄上万字的中文长文本，这种做法会漏掉一个重要的变量：
同样的文字，排在 40 字一行的屏幕上和 20 字一行的书页上，读者付出的努力并不相同。

本文提出一种同时考虑\textbf{语言难度}与\textbf{版面负荷}的中文长文本可读性评估方法。
主要工作如下：
（1）构建了包含 1\,200 篇文章、36\,000 条阅读记录的中文长文本阅读数据集，
记录了每位读者在不同版式下的阅读时间与理解得分；
（2）提出版面负荷指标，用行长、行距与段落密度刻画「排版让眼睛多做了多少事」；
（3）设计了融合语言特征与版面特征的评估模型，在数据集上与人工评分的相关系数达到
$0.81$，比只用语言特征的基线提高 $0.12$。

实验表明，版面因素可以解释可读性差异中约五分之一的部分。
这一结论为中文排版规范、阅读类产品的版式设计提供了可量化的依据。

\vspace{1em}
\noindent{\heiti 关键词：}可读性；中文排版；长文本；阅读负荷；评估模型

\chapter*{Abstract}
\addcontentsline{toc}{chapter}{Abstract}
\markboth{Abstract}{Abstract}

Whether a text is tiring to read depends on what it says and on how it is set.
Classical readability formulas look only at the former. For long-form Chinese text,
which may run to tens of thousands of characters, this ignores an important variable:
the same words cost the reader different effort on a 40-character screen line and on
a 20-character book line.

This thesis proposes a readability assessment method for long-form Chinese text that
accounts for both \emph{linguistic difficulty} and \emph{layout load}. We build a
dataset of 1,200 articles and 36,000 reading sessions, define a layout-load index from
line length, leading and paragraph density, and design a model that combines linguistic
and layout features. The model correlates with human ratings at $r = 0.81$, an
improvement of $0.12$ over a language-only baseline.

\vspace{1em}
\noindent\textbf{Keywords:} readability; Chinese typography; long-form text;
reading load; assessment model

{\hypersetup{linkcolor=black}\tableofcontents}

\mainmatter
\chapter{绪论}
\label{chap:intro}

\section{研究背景}

阅读正在从纸面迁移到屏幕。新闻、小说、技术文档乃至学术论文，
越来越多地以长文本的形式出现在手机、平板和显示器上。
与纸书不同，屏幕上的版式由软件决定：行长随窗口变化，行距由样式表设定，
同一篇文章在不同设备上可能呈现出完全不同的面貌。

可读性研究已有近百年历史\cite{flesch1948}，但经典公式主要面向英文短文，
并且默认版面是固定的。中文没有词间空格，长文本又常常跨越数十屏，
读者的疲劳更多来自「找下一行」「找回刚才读到的位置」，而这些恰恰是版面造成的。

\section{研究问题}

本文试图回答三个问题：
\begin{enumerate}
  \item 在中文长文本中，版面因素对阅读负荷的影响有多大？
  \item 能否用少数几个可计算的指标刻画这种影响？
  \item 把版面指标加入可读性模型后，评估结果能提升多少？
\end{enumerate}

\section{主要贡献}

\begin{enumerate}
  \item 构建了中文长文本阅读数据集，包含 1\,200 篇文章、312 位读者、36\,000 条阅读记录；
  \item 提出版面负荷指标 $\Lambda$，由行长、行距与段落密度三部分组成；
  \item 设计了融合语言与版面特征的评估模型，并开源全部数据与代码。
\end{enumerate}

\section{论文结构}

第~\ref{chap:related}~章回顾相关研究；第~\ref{chap:method}~章介绍版面负荷指标与评估模型；
第~\ref{chap:experiment}~章报告实验结果；第~\ref{chap:conclusion}~章总结全文并讨论后续工作。

\chapter{相关研究}
\label{chap:related}

\section{可读性公式}

Flesch 阅读难度公式\cite{flesch1948}用平均句长和平均音节数估计英文文本的难度，
至今仍被广泛使用。中文方面，研究者以字频、词频和句长为特征，
提出了多种面向教材与新闻的可读性公式\cite{wang2020}。
这些方法的共同点是只看文字本身。

\section{排版与阅读}

Tinker 与 Paterson 在 20 世纪 20 年代系统研究了行长、字号与行距对阅读速度的影响\cite{tinker1929}；
近年的研究把同样的问题搬到了屏幕上\cite{dyson2004}。
中文排版的规范性要求见于《中文排版需求》\cite{w3c2020clreq}，
但它给出的是排版规则，并未量化这些规则对读者的影响。

\section{小结}

已有工作要么只看文字，要么只看版面，
缺少一个把两者放在同一框架下评估的方法。这正是本文要补上的部分。

\chapter{版面负荷与评估模型}
\label{chap:method}

\section{版面负荷指标}

\begin{definition}[版面负荷]
  设一段文本在某一版式下的行长为 $L$（字），行距为 $s$（字号的倍数），
  段落密度为 $d$（每屏段落数），定义其版面负荷为
  \begin{equation}
    \label{eq:load}
    \Lambda(L, s, d) = \alpha \left( \frac{L - L^{*}}{L^{*}} \right)^{2}
      + \beta \, \frac{1}{s} + \gamma \, \frac{1}{1 + d},
  \end{equation}
  其中 $L^{*}$ 为最适行长，$\alpha, \beta, \gamma \geqslant 0$ 为待估参数。
\end{definition}

式~\eqref{eq:load} 的三项分别对应三种额外的眼动成本：
行太长或太短时的回扫与频繁换行、行距过小时的串行，以及段落过长时的定位困难。

\begin{proposition}
  当 $s$ 与 $d$ 固定时，$\Lambda$ 在 $L = L^{*}$ 处取得最小值。
\end{proposition}

\section{评估模型}

模型的输入由语言特征 $\symbf{x}_{\text{lang}}$ 与版面特征 $\Lambda$ 拼接而成，
输出为预测的可读性得分 $\hat{y}$：
\begin{equation}
  \hat{y} = \symbf{w}^{\top} \bigl[ \symbf{x}_{\text{lang}} ; \Lambda \bigr] + b .
\end{equation}
语言特征包括平均句长、常用字比例、专业术语密度等 18 项，详见附录。
参数用最小二乘估计，并以五折交叉验证选择正则化强度。

\chapter{实验与分析}
\label{chap:experiment}

\section{数据集}

我们从新闻、科普、小说与技术文档四类来源中各选取 300 篇文章，
每篇 3\,000～12\,000 字。312 位读者在三种行长、两种行距下阅读，
每读完一篇回答五道理解题，并对疲劳程度打分（1～7）。

\section{整体结果}

表~\ref{tab:main} 比较了三种模型与人工评分的相关系数。
加入版面负荷后，相关系数从 $0.69$ 提高到 $0.81$。

\begin{table}[htbp]
  \centering
  \caption{不同模型与人工评分的相关系数}
  \label{tab:main}
  \begin{tabular}{@{}l c c c@{}}
    \toprule
    模型 & 新闻与科普 & 小说 & 技术文档 \\
    \midrule
    只用语言特征 & 0.71 & 0.66 & 0.70 \\
    只用版面特征 & 0.38 & 0.45 & 0.31 \\
    \textbf{二者融合（本文）} & \textbf{0.82} & \textbf{0.80} & \textbf{0.81} \\
    \bottomrule
  \end{tabular}
\end{table}

\section{行长的影响}

图~\ref{fig:fatigue} 给出了不同行长下读者自评的疲劳程度。
行长在 20～26 字时疲劳最低；超过 34 字后，疲劳随行长近似线性上升。

\begin{figure}[htbp]
  \centering
  \begin{tikzpicture}
    \begin{axis}[
      width=.8\linewidth, height=5.6cm,
      xlabel={行长（字）}, ylabel={疲劳自评（1--7）},
      xmin=12, xmax=46, ymin=2, ymax=5.6,
      axis lines=left, axis line style={mist}, tick style={mist},
      tick label style={font=\small}, label style={font=\small},
      legend style={draw=none, font=\small, at={(0.03,0.97)}, anchor=north west},
    ]
      \addplot[accent, very thick, smooth, mark=*, mark size=1.6pt] coordinates
        {(14,3.6) (18,3.0) (22,2.7) (26,2.8) (30,3.2) (34,3.7) (38,4.3) (42,4.9)};
      \addlegendentry{行距 1.8}
      \addplot[ink, very thick, smooth, dashed, mark=square*, mark size=1.4pt] coordinates
        {(14,3.8) (18,3.3) (22,3.0) (26,3.2) (30,3.7) (34,4.3) (38,4.9) (42,5.4)};
      \addlegendentry{行距 1.4}
    \end{axis}
  \end{tikzpicture}
  \caption{行长与疲劳程度}
  \label{fig:fatigue}
\end{figure}

\section{小结}

版面并非可读性的背景，而是它的一部分。对于技术文档这类信息密集的长文本，
合适的行长与行距带来的改善，几乎与把文字改写得更通俗一样多。

\chapter{总结与展望}
\label{chap:conclusion}

\section{全文总结}

本文从「同样的文字，排法不同，读起来累不累也不同」这一朴素观察出发，
提出了版面负荷指标，并把它与语言特征融合为一个可读性评估模型。
实验表明，版面因素可以解释可读性差异中约五分之一的部分，
在信息密集的技术文档中尤为明显。

\section{不足与展望}

\begin{enumerate}
  \item 实验只覆盖了横排简体中文，竖排与繁体文本有待检验；
  \item 版面负荷指标没有考虑字体本身的差异，后续可以引入字重、字面大小等特征；
  \item 疲劳程度依赖读者自评，将来可以结合眼动数据做更细致的测量。
\end{enumerate}


\backmatter
\printbibliography[heading=bibintoc, title=参考文献]
\chapter{致\quad 谢}

论文写到最后，才明白「可读」二字有多难。

感谢导师何明远教授。五年里，他改过我每一稿的每一处措辞，
却从不替我决定研究的方向；他常说，排版是替读者省力，写论文也是。

感谢实验室的同门，陪我熬过数据集标注最枯燥的那个冬天；
感谢参与阅读实验的 312 位志愿者，你们读完的每一篇文章，都是这篇论文的一部分。

最后，感谢我的父母。你们从不过问我研究的是什么，只问我吃饭了没有。


\end{document}
